% Generated draft: requires mathematical and source review.
\documentclass[11pt]{article}
\usepackage[margin=25mm]{geometry}
\usepackage{fontspec}
\setmainfont{Latin Modern Roman}
\usepackage{amsmath,amssymb,mathtools,unicode-math}
\setmathfont{Latin Modern Math}
\usepackage{enumitem,needspace,xurl,hyperref}
\setlength{\parindent}{0pt}
\setlength{\parskip}{5pt}
\setlength{\emergencystretch}{3em}
\newcommand{\N}{\mathbb{N}}
\newcommand{\Z}{\mathbb{Z}}
\newcommand{\Q}{\mathbb{Q}}
\newcommand{\R}{\mathbb{R}}
\newcommand{\C}{\mathbb{C}}
\newcommand{\F}{\mathbb{F}}
\newcommand{\A}{\mathbb{A}}
\newcommand{\PP}{\mathbb{P}}
\newcommand{\E}{\mathbb{E}}
\newcommand{\eps}{\varepsilon}
\newcommand{\dd}{\,\mathrm{d}}
\begin{document}
\section*{UnsolvedMath ID: 20000369}
\textbf{Source problem code:} catalogue entry
\subsection{Definitions and mathematical statement}
Over \(F=\Q(t)\), consider the homogeneous diagonal quadratic equation
\[
q_t(x)=\sum_{j=1}^s A_j(t)x_j^2=0,\qquad A_j(t)\in\Q(t),
\]
where solubility means a nonzero vector over the relevant field. The source asks for a circle method over \(\Q(t)\), and for a Hasse principle using appropriate local conditions, for example those from discrete valuation rings. Wooley notes difficulty identifying major arcs; the valuation class and analytic setup are not specified.

For linear coefficients \(A_j=c_j+td_j\), put \(q_0=\sum c_jx_j^2\), \(q_1=\sum d_jx_j^2\), and
\[
X=\{q_0=q_1=0\}\subset\PP^{s-1}_{\Q}.
\]
The cited Amer--Brumer result is a solubility equivalence between isotropy of \(q_0+tq_1\) over \(\Q(t)\) and a rational point on \(X\); it is not coefficientwise equality for arbitrary rational-function solutions. This is a quartic del Pezzo surface specifically when \(s=5\) and the intersection is smooth, not for every number of variables. The source asks whether the ``obvious'' local conditions over \(\Q(t)\) capture the Brauer--Manin obstruction for \(X\), and Harari suggests corresponding questions over \(\Q_p(t)\), referring to Harari--Szamuely.

The attached aid treats a smooth complete intersection \(X=\{q_0=q_1=0\}\) over a number field \(k\). For each place \(v\) of \(k\), it reports
\[
X(k_v)\ne\varnothing\quad\Longleftrightarrow\quad
q_0+Tq_1\text{ isotropic over }k_v(T),
\]
and the analogous equivalence over \(k\) and \(k(T)\). If \(X(k_v)\ne\varnothing\) for every \(v\), it additionally reports isotropy over the completion \(k(T)_w\) for every rank-one discrete valuation \(w\) of \(k(T)\), not only the coefficient-field tests. It says closed discriminant points are included through the corank-one first residue form; the full rank/smoothness proof hypotheses are not supplied.

Consequently any Hasse counterexample of this type gives a globally anisotropic pencil passing both all coefficient-field and all rank-one-DVR isotropy tests. If its failure is Brauer--Manin, the tests cannot encode that obstruction merely as Boolean isotropy assertions. For clarity, the Brauer--Manin condition uses adelic points \((P_v)\) and simultaneous reciprocity equations
\[
\sum_v\operatorname{inv}_v\alpha(P_v)=0\quad\text{for every }\alpha\in\operatorname{Br}(X),
\]
which depend on the chosen points across places. The record supplies no explicit Hasse counterexample or full replacement local-global principle, and does not construct the requested general circle method.
\subsection{Short English statement}
Develop local-global and circle-method techniques over \(\Q(t)\), especially for quadratic pencils, and compare their local tests with Brauer--Manin on intersections of two quadrics. The attached aid reports that adelically soluble base intersections pass all rank-one valuation isotropy tests, so such tests alone can miss global failures.
\subsection{Sources}
\begin{itemize}
\item Local UnsolvedMath registry, permanent ID \(20000369\), exact public code \texttt{catalogue entry}; canonical location \texttt{aim-algebraic-number-theory-notes.json}, notes[20].
\item AIM workshop ``Rational and integral points on higher-dimensional varieties,'' May 28, 2014, Problem 1; supplied URL \url{https://aimath.org/pastworkshops/ratlhigherdimvarproblems.pdf}. Dialogue is attributed to Browning, Wooley and Harari. The original invokes Amer--Brumer and work of Harari--Szamuely without full publication details; the literature field is empty. The PDF was not consulted here.
\item Attached research attempt 1, ``All-valuation isotropy of a quadratic pencil and the limits of Boolean local tests,'' without complete bibliographic references or an explicit counterexample.
\end{itemize}
\end{document}
